\documentclass[11pt,a4paper]{article}

\usepackage[utf8]{inputenc}
\usepackage[T1]{fontenc}
\usepackage{lmodern}
\usepackage{amsmath,amssymb,amsthm}
\usepackage{booktabs}
\usepackage{microtype}
\usepackage[margin=2.6cm]{geometry}
\usepackage{listings}
\usepackage{xcolor}
\usepackage[hidelinks]{hyperref}
\usepackage{cleveref}

\lstset{
  basicstyle=\ttfamily\small,
  columns=fullflexible,
  keepspaces=true,
  frame=single,
  framerule=0.3pt,
  rulecolor=\color{gray!50},
  backgroundcolor=\color{gray!5},
  breaklines=true,
}

\theoremstyle{plain}
\newtheorem{theorem}{Theorem}
\newtheorem{proposition}[theorem]{Proposition}
\newtheorem{lemma}[theorem]{Lemma}
\theoremstyle{definition}
\newtheorem{definition}[theorem]{Definition}
\theoremstyle{remark}
\newtheorem{remark}[theorem]{Remark}

\newcommand{\R}{\mathbb{R}}
\newcommand{\Z}{\mathbb{Z}}
\newcommand{\Nat}{\mathbb{N}}

\title{PHVE: a bijective, prefix-searchable spatial code\\
for anatomical volumes}

\author{Paul Guindo\\
\small Altius Academy SNC, Échallens, Vaud, Switzerland\\
\small \texttt{paulguindo@altius-group.ch}}

\date{\today}

\begin{document}
\maketitle

\begin{abstract}
We describe PHVE, a spatial code that assigns to every point of a
standardised anatomical reference volume a short string over a
29-character alphabet. The code is bijective on the voxel grid at
sufficient order, and it is \emph{hierarchical}: truncating a code to its
first $k$ characters yields the code of the enclosing dyadic cell, so a
string prefix denotes exactly a region of space. The construction is the
Hilbert index of Skilling composed with an affine normalisation and a
radix-29 expansion; no novelty is claimed for any of these ingredients.

What the software contributes is the packaging and its verification. We
report an exhaustive collision test on the MNI152 brain mask
($235\,375$ voxels), which fixes the order at which the code becomes
injective, and a benchmark of dyadic-cell range queries showing that a
plain ordered index over an integer column answers them with two binary
searches, independently of the number of results, against a
$k$-d tree that must assemble a scattered set. Every figure in this paper
is produced by a seeded script in the accompanying repository.

We are equally explicit about what the code does \emph{not} do. A companion
experimental campaign refuted several claims made for this construction in
earlier drafts --- bandwidth reduction, a compression gain, and an
inter-patient stability property --- and \Cref{sec:not} states each
refutation with a pointer to the script that produced it. The code is
offered as an indexing and addressing tool, and as nothing more.
\end{abstract}

\noindent\textbf{Keywords:} space-filling curves, Hilbert index, spatial
indexing, medical image addressing, prefix search, reproducible software.

\section{Introduction}

Addressing a location inside a body is awkward. Millimetre coordinates in a
template space are precise but unreadable, carry no notion of scale, and
cannot be compared as strings. Anatomical names are readable but coarse and
ambiguous. Databases that must answer ``what is stored near here'' usually
reach for a spatial index --- an R-tree or a $k$-d tree --- which means an
extension or an auxiliary structure that a plain relational column does not
provide.

Geographic information systems solved the analogous problem long ago with
hierarchical short codes: Geohash, S2, H3, what3words. A location becomes a
string, a shorter string is a coarser region, and containment is a prefix
test, so an ordinary B-tree suffices. PHVE (\emph{Parametric Hilbert Volume
Encoding}) applies the same recipe to anatomical reference volumes.

\paragraph{What is and is not new.} The Hilbert curve, its Skilling
construction \cite{skilling}, and radix encodings of space-filling-curve
indices are classical, and hierarchical geographic codes are widely
deployed. Nothing in \Cref{sec:construction} is new mathematics. The
contribution of this paper is a specific, verified, reproducible
instantiation for anatomy: a reference family of volumes, an alphabet
chosen to avoid transcription errors, an exhaustive determination of the
order at which the code is injective on a standard brain mask, and a
measured statement of what prefix queries cost. We are explicit about this
because an earlier draft of this work claimed more, and much of it did not
survive testing (\Cref{sec:not}).

\section{Construction}\label{sec:construction}

\subsection{Reference family}

A \emph{reference volume} is an axis-aligned box
$V_\alpha=\prod_{i=1}^{3}[a_i,b_i]\subset\R^3$ with an identifier
$\alpha$, fixed once and published. \Cref{tab:volumes} lists the family used
here. Coordinates are millimetres in the template space of the volume; for
the cranial volume \texttt{CR} this is MNI space.

\begin{table}[h]
\centering
\begin{tabular}{llr}
\toprule
Code & Volume & Dimensions (mm)\\
\midrule
\texttt{CR} & Cranium    & $300\times250\times250$\\
\texttt{TX} & Thorax     & $400\times350\times300$\\
\texttt{CA} & Heart      & $150\times120\times100$\\
\texttt{AB} & Abdomen    & $400\times350\times300$\\
\texttt{MB} & Limb       & $700\times200\times200$\\
\texttt{FB} & Full body  & $2000\times600\times400$\\
\bottomrule
\end{tabular}
\caption{The reference family. The identifier is part of the code, so a
code is meaningless without it and cannot be misread across volumes.}
\label{tab:volumes}
\end{table}

\subsection{The map}

Fix an order $p\in\Nat$ and write $n=2^p$. The encoding is the composition
\[
F_{\alpha,p}\;=\;\mathrm{Enc}\circ\mathrm{Hil}_p\circ\mathrm{Nor}_{\alpha,p},
\]
of three maps.

\emph{Normalisation.} $\mathrm{Nor}_{\alpha,p}:V_\alpha\to\{0,\dots,n-1\}^3$
sends $\mathbf{x}$ to the integer cell
\[
g_i \;=\; \min\Bigl(n-1,\;
   \Bigl\lfloor \frac{x_i-a_i}{b_i-a_i}\,n \Bigr\rfloor\Bigr),
\qquad i=1,2,3 .
\]
The cell edge lengths are $\Delta_i=(b_i-a_i)/n$, which differ between axes
when the box is not cubic.

\emph{Hilbert index.} $\mathrm{Hil}_p:\{0,\dots,n-1\}^3\to\{0,\dots,2^{3p}-1\}$
is the Skilling index \cite{skilling}, a bijection between the grid and an
interval of integers.

\emph{Radix expansion.} $\mathrm{Enc}$ writes the integer in base $29$ over
the alphabet
\[
\Sigma=\texttt{23456789ABCDEFGHJKMNPQRSTVWXY},
\]
zero-padded to
$k(p)=\lceil 3p\log 2/\log 29\rceil$ characters. The alphabet omits
\texttt{0}, \texttt{1}, \texttt{I}, \texttt{L}, \texttt{O}, \texttt{U} and
the lower case, so that no two characters are confusable when a code is
read aloud, handwritten, or scanned. At $p=8$ a code is $5$ characters, and
a full address reads
\begin{center}
\texttt{CR:GVP-5X}
\end{center}
for the point $(0,50,30)$\,mm in MNI space, the hyphen being cosmetic.

\subsection{The three properties that matter}

\begin{proposition}[Bijectivity]\label{prop:bij}
For every $p$, $\mathrm{Hil}_p$ is a bijection, hence $F_{\alpha,p}$ is
injective on the set of grid cells and a code determines its cell exactly.
\end{proposition}

This is a property of the Skilling construction, not of ours. What is not
automatic, and is the practical question, is whether distinct \emph{voxels}
of a given image fall in distinct cells; this is settled empirically in
\Cref{sec:collisions}.

\begin{proposition}[Truncation is containment]\label{prop:trunc}
Let $c=F_{\alpha,p}(\mathbf{x})$ and let $c|_k$ denote its first $k$
characters. Then the set of points whose code shares the prefix $c|_k$ is
exactly a dyadic sub-box of $V_\alpha$ containing $\mathbf{x}$, and these
sub-boxes are nested as $k$ decreases.
\end{proposition}

\begin{proof}[Proof sketch]
The Hilbert index has the property that the $2^{3\ell}$ cells of a dyadic
sub-box of level $\ell$ occupy a contiguous, aligned block of indices of
length $2^{3(p-\ell)}$. A prefix in base $29$ fixes the leading digits,
hence an aligned block of integers; the alignment of the two block
structures is what makes the correspondence exact at the levels where
$29^{j}$ and $2^{3\ell}$ align, and a nesting elsewhere.
\end{proof}

\Cref{prop:trunc} is the reason the code is useful: \emph{a region query is
a string prefix}, so it is a range scan on an ordinary index.

\begin{proposition}[Quantisation]\label{prop:quant}
Decoding a code to the centre of its cell yields a point within
$\tfrac12\|\Delta\|_2$ of the original, and this bound is attained.
\end{proposition}

The bound is elementary; that it is \emph{attained} was checked
numerically, the measured maximum error agreeing with the theoretical value
to all printed digits at $p\in\{6,8,10,12\}$.

\begin{remark}[Which Hilbert curve]\label{rem:variant}
There are $10\,694\,807$ structurally distinct three-dimensional Hilbert
curves \cite{haverkort}. This implementation uses Skilling's. In $d=2$ it
coincides exactly with the index of Hamilton and Rau-Chaplin
\cite{hamilton}; in $d=3$ it does \emph{not} --- the two induce different
total orders already at $p=1$, and no composition with any of the $48$
signed axis permutations of the cube (with or without reversal of the
traversal) reconciles them for $p\ge2$. Any locality constant quoted for
this software is therefore a property of the Skilling variant and must not
be compared with published values for other variants. In particular, the
worst-case $L_2$ dilation is $6$ in $d=2$ \cite{moon} but approximately
$29.5$ for the Skilling variant in $d=3$, against $22.9$ for the best
published three-dimensional curve \cite{haverkort}: the widely implemented
variant is about $24\%$ worse. This matters for users, because that
constant is what governs how tightly a prefix confines a point.
\end{remark}

\section{Verification}

All scripts are in \texttt{experiments/} of the accompanying repository,
seeded, and write a JSON record of every number they produce. The machine
used is an AMD Ryzen 7 7840HS (8 cores, L1d 32\,KiB/core, L2 1\,MiB/core,
L3 16\,MiB shared), Linux 6.17, CPython 3.12.3, NumPy 2.4.2, SciPy 1.17.1.

\subsection{Kernel}

The vectorised Skilling kernel is bijective with an exact round trip and
unit-step adjacency for $d\in\{2,3\}$ and $p\le5$, and agrees with the
repository's independent scalar implementation on all $4096$ points at
$d=3$, $p=4$.

\subsection{Injectivity on a real volume}\label{sec:collisions}

The practical question is at which order the code separates the voxels of
an actual image. We encoded the centre of every voxel of the MNI152 brain
mask at 2\,mm (\texttt{nilearn}, $235\,375$ voxels, all inside the
\texttt{CR} box) and counted distinct codes.

\begin{table}[h]
\centering
\begin{tabular}{rlrrr}
\toprule
$p$ & cell (mm) & colliding voxels & rate & max multiplicity\\
\midrule
6  & $4.69\times3.91\times3.91$ & $207\,171$ & $88.02\%$ & 12\\
7  & $2.34\times1.95\times1.95$ & $33\,183$  & $14.10\%$ & 2\\
8  & $1.17\times0.98\times0.98$ & $\mathbf{0}$ & $0.00\%$ & 1\\
9, 10 & --- & $0$ & $0.00\%$ & 1\\
\bottomrule
\end{tabular}
\caption{Collisions on the MNI152 brain mask. Injectivity begins exactly
at $p=8$, which is where the cell first fits inside the 2\,mm voxel along
every axis. Script: \texttt{exp15\_v4\_carryover.py}.}
\label{tab:collisions}
\end{table}

The threshold is not a surprise --- it is where $\Delta_i<2$\,mm for all
$i$ --- but it had never been measured, and $p=8$ is therefore the smallest
order at which a code may be used as a voxel identifier for this template.
At $p=7$ one voxel in seven shares its code with a neighbour.

\subsection{Prefix range queries}\label{sec:bench}

By \Cref{prop:trunc} the points inside a dyadic cell occupy a contiguous
range of indices. On a sorted integer column the query is two binary
searches and a slice; the client supplies the prefix, so no Hilbert
computation happens at query time. We compare against a \texttt{scipy}
$k$-d tree on the same points, using a Chebyshev ball, which \emph{is} the
axis-aligned box, so both methods answer the same question exactly.

We stress that this is a \emph{restricted} query class. A $k$-d tree
answers arbitrary boxes and balls; the prefix index answers dyadic cells
only. Comparing them on arbitrary boxes would be meaningless. The claim
under test is narrow: for the queries the encoding is designed for, no
spatial index is needed.

Before any timing was recorded, both methods were checked to return
identical point sets; over the $540$ queries reported, they always did.

\begin{table}[h]
\centering
\begin{tabular}{rrrrrr}
\toprule
depth & cell (voxels/side) & median results & prefix (ms) & $k$-d tree (ms) & ratio\\
\midrule
2 & 64 & $6849$ & $0.0048$ & $0.4232$ & $87.3$\\
3 & 32 & $1580$ & $0.0034$ & $0.1239$ & $36.3$\\
4 & 16 & $478$  & $0.0030$ & $0.0572$ & $19.0$\\
5 & 8  & $64$   & $0.0028$ & $0.0206$ & $7.5$\\
\bottomrule
\end{tabular}
\caption{Dyadic-cell queries on the $235\,375$-voxel brain mask at $p=8$;
medians over $24$--$200$ queries per depth, all on non-empty cells.
Script: \texttt{exp22\_prefix\_index.py}.}
\label{tab:bench}
\end{table}

The structural point, rather than the ratio, is what should be taken away.
The prefix query cost is essentially independent of how many points it
returns ($0.0028$--$0.0048$\,ms across a hundredfold change in
selectivity), because the answer is a contiguous \emph{range} rather than a
set: locating it costs $O(\log N)$ and reading it is a sequential scan. The
$k$-d tree must assemble scattered indices, so its cost grows with the
output.

\paragraph{Where the prefix index loses.} It is \emph{not} cheaper to
build: $0.123$\,s against $0.056$\,s for the $k$-d tree, since encoding and
sorting is more work than a $k$-d tree construction. It is smaller
($3.6$\,MiB against at least $5.4$\,MiB for the coordinate array alone),
and it requires no spatial extension --- an integer or string column in any
database will do --- but the build cost is a genuine disadvantage and we
report it as such.

\section{What this code does not do}\label{sec:not}

An earlier draft of this work advanced several further claims. A systematic
experimental campaign refuted them. We list them here, with the script that
settled each, because a software paper that omits its own negative results
is not usable by anyone.

\begin{itemize}
\item \textbf{It does not reduce matrix bandwidth.} Renumbering the
vertices of a finite-element mesh by this code leaves the maximum index gap
at $\Theta(N)$, essentially unchanged from the natural ordering, and worse
than reverse Cuthill--McKee by more than an order of magnitude. What it
reduces is the \emph{mean} and \emph{median} gap. Moreover the renumbering
is analytically inert: it is a permutation, so the spectrum, the
conditioning and the solution are unchanged, verified to $10^{-12}$.

\item \textbf{It does not accelerate a sparse solver in any general sense.}
With a permutation-equivariant preconditioner the iteration count is
invariant by construction (exactly so for a Chebyshev polynomial, over all
sizes and seeds tested). The only measured advantages are a $35$--$43\%$
reduction in \emph{simulated} L1 cache misses above $N\approx3\cdot10^4$,
with no effect at L2 or L3, and an ordering that is $1.5$--$4.3\times$
cheaper to compute than reverse Cuthill--McKee. With a fixed-pattern
incomplete Cholesky, which ordering gives fewer iterations was found to
depend on the mesh seed and could not be determined.

\item \textbf{It does not improve DPCM compression} of the volumes tested.
On MNI152 the Hilbert traversal loses to a raster scan. An earlier reported
gain came from a script that quantised a signal on $[0,1]$ with
\texttt{round}, reducing it to two levels before taking residuals.

\item \textbf{It is not the compact Hilbert index} of Hamilton and
Rau-Chaplin in three dimensions (\Cref{rem:variant}), although in two
dimensions it is identical to theirs.

\item \textbf{Its stability under coordinate jitter is quantisation, not
anatomy.} An experiment reported as ``inter-patient stability'' uses no
inter-patient data: it perturbs fixed landmarks with Gaussian noise and
counts unchanged codes. Over $p\in\{5,\dots,8\}$ and
$\sigma\in\{0.25,\dots,4\}$\,mm the measured rate agrees with the
elementary prediction that a point merely has to remain inside its own cell
to within $0.05$, and to within $0.01$ for $\sigma\ge2$\,mm. It is a
property of the grid. Whether codes agree across registered subjects is an
open question that this software does not answer.
\end{itemize}

\section{Software}

\paragraph{Availability.} The reference implementation, the demonstration
application, and every script cited above are available at
\url{https://github.com/Altius-Academy-SNC/PHVE} under the MIT licence.

\paragraph{Interface.} The core is three functions.

\begin{lstlisting}[language=Python]
from codec3d import encode, decode, truncate

code = encode(x_mm, y_mm, z_mm, p=8, volume="CR")   # -> 'CR:GVP-5X'
xyz, resolution_mm = decode(code)                   # cell centre + scale
region = truncate(code, n_chars=3)                  # enclosing dyadic cell
\end{lstlisting}

\paragraph{Use in a database.} Because a region is a prefix, no spatial
extension is required:
\begin{lstlisting}[language=SQL]
CREATE INDEX ON findings (phve_code);
SELECT * FROM findings WHERE phve_code LIKE 'CR:GVP%';
\end{lstlisting}

\paragraph{Demonstration.} A Streamlit application renders the traversal on
a brain volume, the effect of truncation on the located region, and the
prefix query of \Cref{sec:bench}.

\paragraph{Reproducibility.} Every number in this paper comes from a seeded
script that writes a JSON record to \texttt{experiments/results/}. The
laboratory notebook \texttt{experiments/LOGBOOK.md} records the runs in
chronological order, including the failures and the two bugs found in
earlier scripts; \texttt{experiments/UNVERIFIED.md} lists what remains
untested.

\section{Limitations}

The code addresses a \emph{reference} volume. Using it for a subject
requires registering that subject to the template, and the registration
error is then added to the quantisation error of \Cref{prop:quant}; this
paper measures neither. The family of \Cref{tab:volumes} is fixed by
convention and its boxes are not derived from any anatomical criterion.
Injectivity was established for one template at one resolution
(\Cref{tab:collisions}); other templates will have their own threshold,
computable by the same script. Finally, the benchmark of \Cref{sec:bench}
is single-threaded, in-memory, and on one dataset; it says nothing about
behaviour on disk or under concurrency.

\begin{thebibliography}{9}

\bibitem{skilling} J.~Skilling. Programming the Hilbert curve.
\emph{AIP Conference Proceedings} \textbf{707} (2004), 381--387.

\bibitem{hamilton} C.~H. Hamilton and A.~Rau-Chaplin. Compact Hilbert
indices: space-filling curves for domains with unequal side lengths.
\emph{Information Processing Letters} \textbf{105} (2008), 155--163.

\bibitem{haverkort} H.~Haverkort. An inventory of three-dimensional Hilbert
space-filling curves. arXiv:1109.2323 (2011); and
How many three-dimensional Hilbert curves are there? arXiv:1610.00155
(2016).

\bibitem{moon} B.~Moon, H.~V. Jagadish, C.~Faloutsos and J.~H. Saltz.
Analysis of the clustering properties of the Hilbert space-filling curve.
\emph{IEEE Transactions on Knowledge and Data Engineering} \textbf{13}
(2001), 124--141.

\bibitem{gotsman} C.~Gotsman and M.~Lindenbaum. On the metric properties of
discrete space-filling curves. \emph{IEEE Transactions on Image Processing}
\textbf{5} (1996), 794--797.

\end{thebibliography}

\end{document}
